\chapter*{生成AI等の利用に関する申告}
\markboth{生成AI等の利用に関する申告}{生成AI等の利用に関する申告}
\addcontentsline{toc}{chapter}{生成AI等の利用に関する申告}

\UsamiFictionalExampleNotice
\par\vspace{1.5\baselineskip}

本論文の作成にあたり，著者はOpenAI ChatGPT（GPT-5，2026年1月15日から2月20日まで）を，研究課題の論点整理，Python実装の補助，コードのリファクタリング，章構成の検討，文章表現の推敲および批判的確認に使用した．利用範囲は，第4章の前処理コード，第5章の集計コードおよび全章の文章表現である．生成AIの出力は研究上の根拠としてそのまま採用せず，実験設計，ソースコード，モデル出力，数値結果，図表，主張，引用文献および最終稿を，著者が原資料，一次文献および実行ログと照合した．本架空例では，個人情報，認証情報および未公開研究データを外部AIサービスへ入力していない．生成AIを著者として扱わず，本論文の正確性，研究公正，独創性および最終内容については著者が責任を負う．

% 次の引用は申告形式の参考であり，研究上の関係がある場合だけ原論文を読んで用いる．
AI利用の範囲と人間による検証責任を分けて記録する申告形式は，宇佐美ら\cite{usami2026darkcurrent}の開示例を参考にした．
